\documentclass[11pt,oneside]{book}

% --- Engine / fonts ---
\usepackage{fontspec}
\setmainfont{TeX Gyre Pagella}
\setmonofont{DejaVu Sans Mono}[Scale=0.85]
\usepackage{unicode-math}
\setmathfont{Latin Modern Math}

% A handful of unicode math-alphabet characters (blackboard, script, and
% a combining-accent letter) used directly in running prose, outside of
% math mode, are absent from TeX Gyre Pagella; substitute them from
% fallback fonts that cover these ranges.
\usepackage{newunicodechar}
\newfontfamily\symbolfallback{Latin Modern Math}
\newfontfamily\accentfallback{DejaVu Sans}
\newunicodechar{ẑ}{{\accentfallback ẑ}}
\newunicodechar{ℝ}{{\symbolfallback ℝ}}
\newunicodechar{𝒢}{{\symbolfallback 𝒢}}
\newunicodechar{𝒲}{{\symbolfallback 𝒲}}
\newunicodechar{𝓜}{{\symbolfallback 𝓜}}
\newunicodechar{𝒫}{{\symbolfallback 𝒫}}
\newunicodechar{Ṁ}{{\accentfallback Ṁ}}
\newunicodechar{ᵢ}{{\symbolfallback ᵢ}}
\newunicodechar{Č}{{\accentfallback Č}}

% --- Layout ---
\usepackage[a5paper,margin=2.2cm,bindingoffset=0.4cm]{geometry}
\usepackage{setspace}
\onehalfspacing
\usepackage{titlesec}
\titleformat{\chapter}[display]
  {\normalfont\bfseries\Large}{\chaptertitlename\ \thechapter}{10pt}{\Large}
\titlespacing*{\chapter}{0pt}{20pt}{28pt}

% --- Front matter helpers ---
\usepackage{titling}
\usepackage{fancyhdr}
\setlength{\headheight}{26pt}
\pagestyle{fancy}
\fancyhf{}
\fancyhead[LE,RO]{\thepage}
\fancyhead[RE]{\nouppercase{\leftmark}}
\fancyhead[LO]{\nouppercase{\rightmark}}
\renewcommand{\headrulewidth}{0.3pt}

% --- References / TOC ---
\usepackage[unicode]{hyperref}
\hypersetup{
  pdftitle={The Physics of Persistent AI Worlds},
  pdfauthor={Flyxion},
  colorlinks=true,
  linkcolor=black,
  urlcolor=black,
  citecolor=black
}
\usepackage{bookmark}

\usepackage{verbatim}
\usepackage{fvextra}
\fvset{breaklines=true, fontsize=\small}
\usepackage{booktabs}
\usepackage{longtable}
\usepackage{float}

% Loosen justification tolerance slightly given long unhyphenated
% mathematical tokens embedded in running prose.
\tolerance=2000
\emergencystretch=2em

\newcommand{\booksubtitle}{From Token Prediction to World Evolution}
\newcommand{\bookalttitle}{Persistent Generative Worlds: A Field Theory of Semantic Simulation}

\begin{document}

% ===================== TITLE PAGE =====================
\begin{titlepage}
\centering
\vspace*{3cm}
{\Huge\bfseries The Physics of\\[4pt] Persistent AI Worlds\par}
\vspace{1cm}
{\Large \booksubtitle \par}
\vspace{0.6cm}
{\normalsize\itshape \bookalttitle \par}
\vfill
{\Large Flyxion\par}
\vspace{0.3cm}
{\normalsize Independent Researcher\par}
\vspace{1.5cm}
{\normalsize 2026\par}
\vspace*{2cm}
\end{titlepage}

\frontmatter

% ===================== PREFACE (placeholder) =====================
\chapter*{Preface}
\addcontentsline{toc}{chapter}{Preface}
This book is an extended argument for a single claim: that a system generating
observations should be answerable to something that persists independent of
being observed, rather than merely conditioned on whatever happens to be
near at hand. It develops that claim from a first symptom, a stone wall
whose crack moves between two glances, through a complete mathematical and
engineering architecture, into a set of applications and, finally, into
what the architecture implies about identity, minds, and civilization once
taken seriously as more than an engineering convenience.

\chapter*{Roadmap of the Theory}
\addcontentsline{toc}{chapter}{Roadmap of the Theory}
This book develops each major concept through four complementary
perspectives. It begins with intuition, introducing why the concept is
needed. It then develops the architectural role the concept plays within
the semantic field. Once the architecture is established, the mathematical
structures required to formalize it are introduced. Finally, the
engineering implications are explored through concrete implementations.
Not every chapter contains all four perspectives; instead, the book
revisits each concept as progressively richer tools become available.

The book's core cycle recurs throughout, refined further each time it
reappears:
\begin{center}
World \(\to\) Persistent Field (M\textsubscript{t}) \(\to\)
Observation (O\textsubscript{\(\theta\)}) \(\to\) Cues \(\to\)
Affordances \(\to\) Neural Proposal (\(\Delta\)\(\psi\)) \(\to\)
Constraint Projection (\(\Pi\)\textsubscript{C}) \(\to\) Commit \(\to\)
Updated World (M\textsubscript{t+1}) \(\to\) World
\end{center}

\begin{center}
\begin{tabular}{ll}
Part I & Why the cycle is needed \\
Part II & The world and observation \\
Part III & Action and evolution \\
Part IV & Mathematical guarantees \\
Part V & Software realization \\
Part VI & Applications \\
Part VII & Philosophical consequences \\
\end{tabular}
\end{center}

\tableofcontents

\mainmatter

% ===================== PART I =====================
\part{Why AI Forgets}
\input{chapters/ch01.tex}
\input{chapters/ch02.tex}
\input{chapters/ch03.tex}
\input{chapters/ch04.tex}
\input{chapters/ch05.tex}

% ===================== PART II =====================
\part{The Semantic Field}
\input{chapters/ch06.tex}
\input{chapters/ch07.tex}
\input{chapters/ch08.tex}
\input{chapters/ch09.tex}
\input{chapters/ch10.tex}
\input{chapters/ch11.tex}

% ===================== PART III =====================
\part{Dynamics}
\input{chapters/ch12.tex}
\input{chapters/ch13.tex}
\input{chapters/ch14.tex}
\input{chapters/ch15.tex}
\input{chapters/ch16.tex}
\input{chapters/ch17.tex}

% ===================== PART IV =====================
\part{The Mathematics of Persistence}
\input{chapters/ch18.tex}
\input{chapters/ch19.tex}
\input{chapters/ch20.tex}
\input{chapters/ch21.tex}
\input{chapters/ch22.tex}
\input{chapters/ch23.tex}
\input{chapters/ch24.tex}

% ===================== PART V =====================
\part{Engineering the World Engine}
\input{chapters/ch25.tex}
\input{chapters/ch26.tex}
\input{chapters/ch27.tex}
\input{chapters/ch28.tex}
\input{chapters/ch29.tex}
\input{chapters/ch30.tex}

% ===================== PART VI =====================
\part{Applications}
\input{chapters/ch31.tex}
\input{chapters/ch32.tex}
\input{chapters/ch33.tex}
\input{chapters/ch34.tex}
\input{chapters/ch35.tex}
\input{chapters/ch36.tex}

% ===================== PART VII =====================
\part{Toward Semantic Physics}
\input{chapters/ch37.tex}
\input{chapters/ch38.tex}
\input{chapters/ch39.tex}
\input{chapters/ch40.tex}
\input{chapters/ch41.tex}
\input{chapters/ch42.tex}

% ===================== APPENDICES =====================
\appendix
\input{chapters/apxA.tex}
\input{chapters/apxB.tex}
\input{chapters/apxC.tex}
\input{chapters/apxD.tex}
\input{chapters/apxE.tex}
\input{chapters/apxF.tex}

\end{document}
